\documentclass[11pt,letterpaper,twoside]{article}
\usepackage[english]{babel}
\usepackage{amssymb,amsmath}
\usepackage{fancyhdr}
\usepackage{graphicx}
\usepackage{wrapfig}

 \oddsidemargin  0in \evensidemargin 0in
 \topmargin   -0.25in \headheight 0.25in \headsep 0.25in
 \textwidth   6.5in \textheight 8.75in \marginparsep 0pt \marginparwidth 0pt
 \parskip 1ex  \parindent 0ex \footskip 20pt

\newfont{\bssten}{cmssbx10}
\newfont{\bssnine}{cmssbx10 scaled 900}

%%%%%%%%%%%%%%%%%%%%%%%
\newcommand{\whatizit}{Lecture 14: The Binomial Distribution}
%%%%%%%%%%%%%%%%%%%%%%%


\pagestyle{fancy}  
\fancyhead{\bssnine STOR 155,  \whatizit}
\fancyhead[RE]{} \fancyhead[LO]{}
\fancyhead[LE]{\bssnine \thepage} \fancyhead[RO]{\bssnine \thepage}
\lfoot{} \cfoot{} \rfoot{}   


\newcommand{\var}{\mathrm{Var}}



\begin{document}



\thispagestyle{empty} \vspace*{-0.75in}

{\bssten STOR 155: Introduction to Data Models and Inference \hfill June 6, 2024 \\
Prof. Will Lassiter  \hfill Page 1 of \pageref{totalpag}}
\vspace{10pt}
\begin{center} {{\Large \bf \whatizit}} \end{center}

Running example: the Milgram experiment \vspace{200pt}


{\bf The Binomial Setting} \vspace{200pt}

Examples: \vspace{100pt}

Non-example:

\newpage

{\bf The Binomial Distribution} \vspace{6pt}

Two common use cases: \vspace{120pt}

Possible values for $X$: \vspace{10pt}

Probabilities: \vspace{200pt}

Example: You play Monopoly with your siblings a lot --- so much so that you have found that you win roughly 1 out of every 20 games played.  If you and your siblings will play 20 (independent) games of Monopoly this week, what's the probability that ...

\begin{itemize}

\item ... you win at least once? \vspace{30pt}

\item ... you win exactly twice? \vspace{30pt}

\item ... you win at least twice? \vspace{30pt}

\item ... you win between 1 and 5 games?

\end{itemize}

\newpage

Example: According to a 2020 Gallup poll, 43\% of American voters identify with the Mario Party.  Suppose I sample $n$ voters at random and ask for their party identification. If $X$ represents the number of sampled voters who identify with the Mario Party, then $X$ has a \_\_\_\_\_\_\_\_\_\_\_\_\_\_ distribution.

\begin{itemize}

\item If I sample 300 American voters, how many would I expect to identify with the Mario Party? \vspace{30pt}

\item If I sample $n$ American voters, how many would I expect to identify with the Mario Party? \vspace{30pt}

\end{itemize}

Mean and variance for a binomial random variable \vspace{150pt}

Back to Milgram example: \vspace{80pt}

{\bf Connection between binomial and normal distributions}

\newpage

Example:  In 100 tosses of a fair coin ... 

\begin{itemize}

\item ... what's the probability of 35 or fewer heads? \vspace{80pt}

\item ... what's the probability that the number of heads is within 10\% of the expected number? \vspace{80pt}

\end{itemize}

Example: A college has offered admission to 1200 students this cycle. The number of admitted students who will accept the offer is random; from past experience, roughly 75\% of admitted students accept. 

\begin{itemize}

\item Mean and standard deviation of the number accepting admission? \vspace{80pt}

\item Probability that more than 950 students will accept? \vspace{80pt}

\item If they decide to admit an additional 100 students off the wait list, what's the probability that more than 950 students in total will accept?

\end{itemize}

\label{totalpag}
\end{document}
